\documentclass[a4paper]{article}
\usepackage[utf8]{inputenc}
\usepackage{amsmath,amsfonts,amssymb}
\usepackage{amsthm}

\theoremstyle{plain}
\newtheorem{theorem}{Theorem}
\newtheorem{proposition}[theorem]{Proposition}
\newtheorem{corollary}[theorem]{Corollary}
\newtheorem{lemma}[theorem]{Lemma}
\newtheorem{claim}[theorem]{Claim}
\newtheorem{conjecture}[theorem]{Conjecture}
\numberwithin{theorem}{section}
\renewcommand\thetheorem{\!\!}

\theoremstyle{definition}
\newtheorem{definition}[theorem]{Definition}
\newtheorem{remark}[theorem]{Remark}
\newtheorem{notation}[theorem]{Notation}

\newcommand{\sh}[2]{\mathrm{sh}_{#1}(#2)}
\newcommand{\cyc}[2]{\mathrm{cyc}_{#1}(#2)}
\newcommand{\Ccyc}[2]{\mathrm{Cyc}_{#1}(#2)}
\newcommand{\cF}{\mathcal{F}}

\renewcommand{\Pr}{\mathbb{P}}

\newcommand{\ba}{{\mathbf{a}}}
\newcommand{\bb}{{\mathbf{b}}}
\newcommand{\be}{{\mathbf{e}}}
\newcommand{\hE}{{\hat{E}}}
\newcommand{\hF}{{\hat{F}}}
\newcommand{\tO}{{\widetilde{O}}}

\begin{document}

\section*{Definitions-Only Skeleton}

\subsection*{Elementary Notation Conventions}

Throughout, $n$ is a positive integer.

The notation $[n]$ denotes the finite set $\{1,2,\ldots,n\}$.

A uniform random function $f:[n]\to[n]$ means a function chosen uniformly from the finite set of all functions from $[n]$ to $[n]$.

\subsection*{Source Definitions}

Let $f:[n]\to[n]$, and let $G_f$ be its associated directed graph, with vertex set $[n]$ and directed edges $(i,f(i))$ for $i \in [n]$. A vertex is \emph{cyclic} if it belongs to a cycle of $G_f$.

\subsection*{Target Statement}

\begin{theorem}
	The probability that a uniform random function $f:[n]\to [n]$ has a unique cyclic vertex is~$\frac{1}{n}$.
\end{theorem}

\subsection*{Source Statements Not Granted as Support}

The following source statements are retained only as non-proof mathematical statements from the paper. They are \emph{excluded} under the run's definitions-only support policy.

\begin{itemize}
	\item Functions $f$ having a unique cyclic vertex are in bijection with rooted trees on $[n]$ (the unique cyclic vertex is the root, and edges are oriented towards the root).
	\item Cayley's formula asserts that there are $n^{n-2}$ trees on the vertex set $[n]$, or equivalently there are $n^{n-1}$ rooted trees.
	\item Let $H_n$ denote the height of a uniform random vertex in a uniform random rooted tree with $n$ vertices. Then $1+H_n$ has the same law as the time of first repetition in a i.i.d. sequence of uniform variables on~$[n]$ (time of first collision in the coupon collector process).
\end{itemize}

\end{document}
